\chapter{Programming model}

This chapter describes the rules, patterns, and primitives for writing kernel
code that respects the async-first, shard-per-core model. These invariants are
derived from sitas and are enforced by the kernel's own concurrency design.

\section{Invariants}

\subsection{Only the owning LP mutates its state}

Per-LP data is accessed only from the LP that owns it. Cross-LP mutation must
go through message dispatch (\texttt{ShardMailbox<M>}, IPI closure, or IPI
RPC), never through a shared lock acquired from a foreign LP.

\subsection{Messages are owned values}

Everything that crosses a shard boundary is an owned value --- a
\texttt{Box<dyn FnOnce>}, an \texttt{IpiRpc}, or a typed \texttt{M} in a
\texttt{ShardSender<M>}. No shared references, no \texttt{Arc<Mutex<...>>} as
the normal programming model. \texttt{Send + 'static} boundaries enforce this
at compile time.

\subsection{References to shard-local state must not escape}

\texttt{ShardLocal<T>::with(fn)} provides \texttt{\&mut T} inside a closure;
the borrow is verified at runtime (owner-check + re-entrancy flag). The closure
must not store the reference, send it to another thread, or cross an
\texttt{.await}. The same rule applies to sitas's \texttt{ShardLocal<T>}
userspace equivalent.

\subsection{Work stays on its assigned LP}

A thread spawned on LP $N$ stays on LP $N$. To submit work to another LP, use
\texttt{try\_run\_on\_lp(target, closure)} or \texttt{ShardSender::try\_send}.
The kernel's IPI infrastructure and the sitas \texttt{ShardedSubmitter} are
the two sides of this contract.

\subsection{Observability returns owned snapshots}

The kernel's house style is owned values for diagnostics. sitas's
\texttt{RuntimeSnapshot} / \texttt{ShardSnapshot} pattern applies: never
expose borrowed references into runtime internals. Self-tests and the demo
illustrate this.

\section{Primitives}

\subsection{ShardLocal<T>} Lock-free per-LP container. Use for single-LP,
thread-only data. \\Access: \texttt{sl.with(|val| \{ *val = 42; \})}.

\subsection{ShardMailbox<M>} Typed bounded per-LP queue. Cloneable
\texttt{ShardSender<M>} (any LP can send), single-consumer\\
\texttt{ShardReceiver<M>} (one per LP). \texttt{try\_send(m)} returns
\texttt{Err(m)} on backpressure.

\subsection{IPI closure} \texttt{try\_run\_on\_lp(target\_lp, || \{ ... \})}.
Boxes arbitrary work and delivers it to the target LP's IPI queue. Returns the
closure on full queue.

\subsection{Completion API} \texttt{submit}, \texttt{submit\_worker},
\texttt{complete}, \texttt{poll}, \texttt{wait}, \texttt{cancel}, \texttt{close}.
The exit-observer path \\ (\texttt{submit\_worker}) is the intended default; the
explicit \texttt{complete} is the escape hatch for non-thread-based work.

\subsection{CQ ring} \texttt{CompletionQueueRing} for zero-syscall completion
draining. Kernel side: \texttt{write(cap, result)}. Consumer side:
\texttt{read() -> Option<CqEntry>}. Integration: \texttt{open\_as\_with\_cq}
attaches a ring to a capability table; \texttt{complete()} writes to it.

\section{Rule-of-thumb guidance}

\begin{itemize}
    \item \textbf{Use} \texttt{PerLp<T>} for interrupt-accessed data and
        cross-LP access; \textbf{use} \texttt{ShardLocal<T>} for single-LP,
        thread-only data.
    \item \textbf{Use} \texttt{ShardMailbox<M>} (or \texttt{try\_run\_on\_lp})
        for cross-LP communication; do \textbf{not} lock an LP's state from
        another LP.
    \item \textbf{Use} \texttt{submit\_worker} for fire-and-forget async
        kernel work; the worker just returns, and the exit-observer completes
        the capability.
    \item \textbf{Register} a CQ ring when creating an address space so that
        completions are visible to userspace without per-entry syscalls.
    \item \textbf{Accept} backpressure: bounded tables and queues return
        \texttt{WouldBlock} / \texttt{Err(m)}; retry or apply backpressure
        upstream. Never silently drop messages.
\end{itemize}

\endinput
